\chapter{Spectral Bridges Among Supersymmetric Yang--Mills Cousins}
\label{ch:susy-yang-mills-family-spectra}

\ConventionsLorentzian

The preceding chapters developed three separate structures: pure
\(\mathcal N=1\) Yang--Mills chiral dynamics, pure \(\mathcal N=2\)
Seiberg--Witten special geometry, and planar \(\mathcal N=4\) spectral
integrability.  They are connected by relevant deformations, but the
connection has several layers.  Holomorphic quantities can often be matched
by scale coordinates and protected Ward identities.  Massive spectra,
confining strings, and line-operator Hilbert spaces require additional
nonperturbative data.  This chapter defines the comparison problem.

\section{Comparison Datum}
\label{sec:susy-ym-family-comparison-datum}

Let \(G=SU(N_c)\) with \(N_c\ge2\).  A comparison among the
\(\mathcal N=4,\mathcal N=2,\mathcal N=1\), and bosonic Yang--Mills theories
is a comparison among four QFT problems, not a single family of Lagrangian
symbols.  The amount of supersymmetry determines which quantities are
protected, which deformations preserve a chiral coordinate chart, and which
spectral questions remain dynamical.

\begin{definition}[Yang--Mills cousin comparison datum]
\label{def:ym-cousin-comparison-datum}
A Yang--Mills cousin comparison datum consists of:
\begin{enumerate}[label=(\roman*)]
\item a gauge group \(G=SU(N_c)\) or a specified global form with a specified
  line-operator lattice;
\item a UV description of the \(\mathcal N=4\) theory in the
  \(\mathcal N=1\) field variables
  \[
    V,\qquad \Phi_1,\Phi_2,\Phi_3\in\operatorname{ad}G,
  \]
  with holomorphic coupling
  \[
    \tau=\frac{\theta}{2\pi}+\frac{4\pi\ii}{g_{\rm YM}^2}
  \]
  in the trace convention fixed in Chapter~\ref{ch:susy-gauge-theory};
\item a finite list of relevant deformation parameters, for example adjoint
  chiral masses \(m_i\), an \(\mathcal N=2\to\mathcal N=1\) mass coordinate
  \(\mu\), and a soft gaugino-mass spurion \(m_\lambda\);
\item for each deformation regime, a declared set of observables:
  gauge-invariant local operators, center-charged line operators, domain-wall
  boundary conditions, and defect operators on Wilson lines;
\item a status label for each comparison: holomorphic scale matching,
  protected \(F\)-term matching, semiclassical Abelian approximation, lattice
  spectral datum, planar-integrability datum, or open nonperturbative
  construction.
\end{enumerate}
Two rows of the datum are compared only after the observable type and status
label have been fixed.
\end{definition}

The basic ladder is:
\[
\begin{array}{ccl}
\mathcal N=4
&\xrightarrow{\;m_{\rm hyper}\;}
&\mathcal N=2^\ast
\xrightarrow{\;E\ll |m_{\rm hyper}|\;}
\text{pure }\mathcal N=2,\\[3pt]
\mathcal N=4
&\xrightarrow{\;m_1,m_2,m_3\;}
&\mathcal N=1^\ast
\xrightarrow{\;E\ll \min_i|m_i|\;}
\text{pure }\mathcal N=1,\\[3pt]
\text{pure }\mathcal N=2
&\xrightarrow{\;\mu\,u\;}
&\text{softly broken }\mathcal N=1\text{ with Abelian confinement},\\[3pt]
\text{pure }\mathcal N=1
&\xrightarrow{\;X_\lambda S+\mathrm{c.c.}\;}
&\text{bosonic Yang--Mills deformation problem}.
\end{array}
\]
The arrows are deformation regimes.  A continuum construction for the
endpoint theory, or a proof that a path of deformations reaches a chosen
endpoint without crossing a phase boundary in the observables of interest, is
extra information.

\section{Holomorphic Scale Matching}
\label{sec:susy-ym-family-scale-matching}

Holomorphic scale matching is the cleanest part of the comparison.  It is a
statement about chiral source coordinates and threshold matching in a
holomorphic Wilsonian scheme.  It does not compute the nonchiral spectrum.

\begin{definition}[Holomorphic threshold matching]
\label{def:holomorphic-threshold-matching}
Let a four-dimensional supersymmetric gauge theory with gauge algebra
\(\mathfrak{su}(N_c)\) have holomorphic one-loop coefficient \(b_{\rm high}\)
above a chiral mass threshold \(m\) and coefficient \(b_{\rm low}\) below the
threshold.  A holomorphic threshold matching at the scale \(m\) is the
identity of chiral scale coordinates
\[
  \Lambda_{\rm low}^{b_{\rm low}}
  =
  m^{\,b_{\rm low}-b_{\rm high}}\,
  \Lambda_{\rm high}^{b_{\rm high}},
  \label{eq:holomorphic-threshold-matching}
\]
up to finite holomorphic scheme rescalings already included in the
definitions of \(m\) and \(\Lambda\).
\end{definition}

\begin{proposition}[\(\mathcal N=4\) to pure \(\mathcal N=2\) matching]
\label{prop:n4-to-n2-star-scale-matching}
Write \(\mathcal N=4\) Yang--Mills in \(\mathcal N=2\) variables as one
\(\mathcal N=2\) vector multiplet plus one adjoint hypermultiplet.  Give the
hypermultiplet a holomorphic mass \(m_{\rm h}\).  In the region
\[
  E\ll |m_{\rm h}|,
\]
the remaining pure \(\mathcal N=2\) holomorphic scale satisfies
\[
  \Lambda_{\mathcal N=2}^{2N_c}
  =
  m_{\rm h}^{2N_c}\,q,
  \qquad
  q=\ee^{2\pi\ii\tau}.
  \label{eq:n4-to-n2-scale-matching}
\]
\end{proposition}

\begin{proof}
The \(\mathcal N=4\) theory has vanishing holomorphic beta coefficient, so
the high-energy scale coordinate is the dimensionless instanton factor
\(q=\ee^{2\pi\ii\tau}\).  Pure \(\mathcal N=2\) \(SU(N_c)\) Yang--Mills has
holomorphic coefficient
\[
  b_{\mathcal N=2}=2N_c .
\]
Applying \eqref{eq:holomorphic-threshold-matching} with
\[
  b_{\rm high}=0,\qquad b_{\rm low}=2N_c,\qquad m=m_{\rm h}
\]
gives \eqref{eq:n4-to-n2-scale-matching}.  The dimension of the right-hand
side is \(2N_c\), matching the dimension of
\(\Lambda_{\mathcal N=2}^{2N_c}\).
\end{proof}

\begin{proposition}[\(\mathcal N=4\) to pure \(\mathcal N=1\) matching]
\label{prop:n4-to-n1-star-scale-matching}
Write \(\mathcal N=4\) Yang--Mills in \(\mathcal N=1\) variables as one
vector multiplet and three adjoint chiral multiplets
\(\Phi_1,\Phi_2,\Phi_3\).  Give the three adjoint chiral multiplets
holomorphic masses \(m_1,m_2,m_3\).  In the region
\[
  E\ll \min_i |m_i|,
\]
the pure \(\mathcal N=1\) holomorphic scale satisfies
\[
  \Lambda_{\mathcal N=1}^{3N_c}
  =
  (m_1m_2m_3)^{N_c}\,q,
  \qquad
  q=\ee^{2\pi\ii\tau}.
  \label{eq:n4-to-n1-scale-matching}
\]
\end{proposition}

\begin{proof}
The high-energy \(\mathcal N=4\) coefficient is zero.  Pure
\(\mathcal N=1\) \(SU(N_c)\) Yang--Mills has coefficient
\[
  b_{\mathcal N=1}=3N_c .
\]
Each massive adjoint chiral multiplet changes the holomorphic coefficient by
\(T(\operatorname{ad})=N_c\).  Applying
\eqref{eq:holomorphic-threshold-matching} successively at the three mass
thresholds gives
\[
  \Lambda_{\mathcal N=1}^{3N_c}
  =
  m_1^{N_c}m_2^{N_c}m_3^{N_c}q,
\]
which is \eqref{eq:n4-to-n1-scale-matching}.  The mass dimension is
\(3N_c\), as required for the pure \(\mathcal N=1\) scale coordinate.
\end{proof}

\begin{example}[\(\mathcal N=1^\ast\) chiral critical equations]
\label{ex:n1-star-critical-equations}
In an \(\mathcal N=1\) presentation, the equal-mass \(\mathcal N=1^\ast\)
superpotential is
\[
  W
  =
  \sqrt2\,\operatorname{tr}\Phi_1[\Phi_2,\Phi_3]
  +
  \frac{m}{2}\sum_{i=1}^3\operatorname{tr}\Phi_i^2 .
  \label{eq:n1-star-superpotential}
\]
The \(F\)-term equations are
\[
  \sqrt2[\Phi_2,\Phi_3]+m\Phi_1=0,
  \qquad
  \sqrt2[\Phi_3,\Phi_1]+m\Phi_2=0,
  \qquad
  \sqrt2[\Phi_1,\Phi_2]+m\Phi_3=0.
  \label{eq:n1-star-f-terms}
\]
If \(J_i\) are matrices obeying
\[
  [J_i,J_j]=\ii\epsilon_{ijk}J_k,
  \label{eq:su2-fuzzy-sphere-algebra}
\]
then
\[
  \Phi_i=\frac{\ii m}{\sqrt2}\,J_i
  \label{eq:n1-star-fuzzy-sphere-solution}
\]
solves \eqref{eq:n1-star-f-terms}.  Such solutions organize chiral vacua of
the massive adjoint theory by embeddings of \(\mathfrak{su}(2)\) into
\(\mathfrak{su}(N_c)\).  The full massive spectrum and confinement data also
depend on the nonchiral dynamics and on the chosen line-operator lattice.
\end{example}

\section{What Is Uniform Across the Ladder}
\label{sec:susy-ym-family-uniform-data}

The uniform part of the ladder is a map of observables, not a universal
numerical formula.  Four classes appear repeatedly.

\begin{definition}[Family observable ledger]
\label{def:susy-ym-family-observable-ledger}
For a Yang--Mills cousin comparison datum, the observable ledger is the
following table.
\begin{center}
\small
\begin{tabular}{p{0.22\textwidth}|p{0.39\textwidth}|p{0.28\textwidth}}
\textbf{observable type} & \textbf{representative datum} & \textbf{status}\\
\hline
chiral local &
  \(S,\ u,\ \operatorname{tr}\Phi^k\) &
  holomorphic or protected\\
massive local spectral &
  \(\rho_{ij}(B)\) for gauge-invariant sources &
  Hilbert-space or lattice spectral datum\\
center-charged line &
  \(W_R(C),\ H_m(C)\), flux-sector transfer matrix &
  line-operator datum\\
defect spectral &
  \(\mathcal H_{\rm def},\ H_{\rm def},\mathbb O_a(t)\) &
  Wilson-line or cusp datum
\end{tabular}
\end{center}
where \(S\) is the pure \(\mathcal N=1\) glueball coordinate of
Section~\ref{sec:pure-sym-gaugino-condensation}, \(u\) is the rank-one
\(\mathcal N=2\) Coulomb coordinate of
Chapter~\ref{ch:susy-four-dimensional-n2-seiberg-witten}, \(\rho_{ij}\) is
the spectral measure of
Definition~\ref{def:pure-sym-vacuum-spectral-datum}, \(W_R(C)\) is a Wilson
line in representation \(R\), \(H_m(C)\) is a magnetic line with magnetic
charge \(m\), and \(\mathcal H_{\rm def}\) is the defect Hilbert space of
Definition~\ref{def:planar-n4-wilson-line-flux-tube-datum}.
\end{definition}

\begin{proposition}[Observable-type preservation under a controlled threshold]
\label{prop:observable-type-threshold-preservation}
Consider a relevant mass deformation with threshold \(M\) between two
supersymmetric Yang--Mills cousins.  Assume:
\begin{enumerate}[label=(\roman*)]
\item a Wilsonian description exists above and below \(M\) in a common
  regulator or in compatible holomorphic schemes;
\item the local source under consideration has external momenta
  \(|p_a|\ll |M|\);
\item every massive field integrated out has a gap bounded below by
  \(c|M|\) for some \(c>0\) in the chosen background;
\item the line-operator sector is kept fixed by declaring the global form and
  allowed electric and magnetic charges before the deformation.
\end{enumerate}
Then chiral local sources match by holomorphic threshold coordinates, local
nonchiral correlators match after adding local higher-dimension Wilsonian
operators, and center-charge labels of line operators are preserved as
labels.  The numerical masses, string tensions, and line-defect excitation
energies are dynamical functions of the deformation parameters.
\end{proposition}

\begin{proof}
For chiral local sources, the supersymmetric Wilsonian action separates the
holomorphic \(F\)-terms from \(D\)-terms.  The threshold contribution of a
massive chiral multiplet to the holomorphic gauge coupling is one-loop exact
in the Wilsonian holomorphic coordinate, giving the scale matching of
Definition~\ref{def:holomorphic-threshold-matching}.  Thus the chiral
coordinates match by finite holomorphic maps such as
\eqref{eq:n4-to-n2-scale-matching} and \eqref{eq:n4-to-n1-scale-matching}.

For nonchiral local correlators at momenta \(|p_a|\ll |M|\), the massive
propagators admit an expansion in powers of \(p/M\) after the regulator and
subtraction scheme have been fixed.  The coefficients are local functionals
of the light fields and sources, hence enter the Wilsonian action as
higher-dimension local operators.  This is an equality of low-momentum
Wilsonian coordinates at fixed order in the derivative expansion, with a
remainder controlled only when the chosen regulator supplies the corresponding
analytic estimates.

Line operators require the global form and charge lattice as part of the
definition.  A mass threshold changes the dynamics of screening and
confinement, but it does not change the declared lattice of allowed line
charges.  Therefore the center or charge label can be transported through the
comparison datum.  The area-law coefficient or defect excitation energy is
obtained from the Hilbert-space or transfer-matrix problem in the deformed
theory, so it remains a dynamical spectral quantity.
\end{proof}

\section[Center Strings and k-String Ledgers]{Center Strings and \(k\)-String Ledgers}
\label{sec:susy-ym-family-k-string-ledgers}

For \(SU(N_c)\) with the usual Wilson lines, the center charge of a
representation is an element of \(\mathbb Z_{N_c}\).  A confining theory may
have stable flux sectors labeled by
\[
  k\in\{0,1,\ldots,N_c-1\},
\]
with \(k\) and \(N_c-k\) related by charge conjugation.

\begin{definition}[\(k\)-string tension datum]
\label{def:k-string-tension-datum}
A \(k\)-string tension datum for a confining Yang--Mills theory consists of:
\begin{enumerate}[label=(\roman*)]
\item a regulator or continuum Hilbert-space construction with spatial
  transfer matrix \(T_L\) on a box whose long direction has length \(L\);
\item a center-flux sector \(k\in\mathbb Z_{N_c}\);
\item the lowest energy \(E_k(L)\) in that sector after the vacuum energy is
  subtracted;
\item a limit
  \[
    T_k=\lim_{L\to\infty}\frac{E_k(L)}{L}
    \label{eq:k-string-transfer-tension}
  \]
  whenever the limit exists.
\end{enumerate}
Equivalently, in an area-law regime for a large rectangular Wilson loop in a
representation of \(N\)-ality \(k\),
\[
  \langle W_k(C)\rangle
  =
  \exp\!\bigl(-T_k\,\operatorname{Area}(C)+o(\operatorname{Area}(C))\bigr)
  \label{eq:k-string-area-law}
\]
defines the same tension if the transfer-matrix and Wilson-loop formulations
have been proved equivalent for the regulator under consideration.
\end{definition}

Two comparison formulae are useful as ledgers:
\[
  R_k^{\rm sine}(N_c)
  =
  \frac{\sin(\pi k/N_c)}{\sin(\pi/N_c)},
  \qquad
  R_k^{\rm Casimir}(N_c)
  =
  \frac{k(N_c-k)}{N_c-1}.
  \label{eq:k-string-ledger-ratios}
\]
The first appears in Abelianized supersymmetric regimes and in several
string-inspired large-\(N_c\) calculations.  The second is a quadratic
color-factor comparison coordinate.  In the monograph they are ledgers for
comparison unless a chapter supplies a derivation in a specified QFT
framework.

\begin{proposition}[Finite algebra of the two \(k\)-string ledgers]
\label{prop:k-string-ledger-algebra}
For \(1\le k\le N_c-1\),
\[
  R_k^{\rm sine}(N_c)=R_{N_c-k}^{\rm sine}(N_c),
  \qquad
  R_k^{\rm Casimir}(N_c)=R_{N_c-k}^{\rm Casimir}(N_c).
  \label{eq:k-string-charge-conjugation-symmetry}
\]
At fixed \(k\) and large \(N_c\),
\[
  R_k^{\rm sine}(N_c)
  =
  k-\frac{\pi^2}{6}\frac{k(k^2-1)}{N_c^2}
  +O(N_c^{-4}),
  \label{eq:sine-law-large-n-expansion}
\]
while
\[
  R_k^{\rm Casimir}(N_c)
  =
  k-\frac{k(k-1)}{N_c}+O(N_c^{-2}).
  \label{eq:casimir-large-n-expansion}
\]
\end{proposition}

\begin{proof}
The symmetry follows from
\[
  \sin\!\left(\frac{\pi(N_c-k)}{N_c}\right)
  =
  \sin\!\left(\pi-\frac{\pi k}{N_c}\right)
  =
  \sin\!\left(\frac{\pi k}{N_c}\right)
\]
and from the equality
\[
  (N_c-k)k=k(N_c-k).
\]
For the sine ledger, expand
\[
  \sin\frac{\pi k}{N_c}
  =
  \frac{\pi k}{N_c}
  -
  \frac{\pi^3k^3}{6N_c^3}
  +O(N_c^{-5}),
  \qquad
  \sin\frac{\pi}{N_c}
  =
  \frac{\pi}{N_c}
  -
  \frac{\pi^3}{6N_c^3}
  +O(N_c^{-5}).
\]
Dividing gives
\[
  R_k^{\rm sine}
  =
  k\left(1-\frac{\pi^2k^2}{6N_c^2}\right)
  \left(1+\frac{\pi^2}{6N_c^2}\right)
  +O(N_c^{-4}),
\]
which is \eqref{eq:sine-law-large-n-expansion}.  For the Casimir ledger,
\[
  \frac{k(N_c-k)}{N_c-1}
  =
  k\frac{1-k/N_c}{1-1/N_c}
  =
  k\left(1-\frac{k}{N_c}\right)
  \left(1+\frac{1}{N_c}+O(N_c^{-2})\right),
\]
which gives \eqref{eq:casimir-large-n-expansion}.
\end{proof}

\section{Entries of the Family}
\label{sec:susy-ym-family-entries}

\subsection{Pure Bosonic Yang--Mills}

Pure bosonic Yang--Mills is the endpoint with the fewest protected
observables and the strongest need for a nonperturbative construction.  Its
glueball and flux-tube spectra are formulated in
Section~\ref{sec:lattice-ym-glueballs-flux-observables}.  In the family
ledger, the bosonic row contains:
\[
\begin{array}{c|c}
\text{datum} & \text{role}\\
\hline
\text{OS-positive lattice gauge measure} &
  \text{candidate regulator and transfer matrix}\\
\text{local operators }\operatorname{tr}F^2,\operatorname{tr}F\widetilde F,\ldots &
  J^{PC}\text{ glueball sources}\\
\text{center-flux sectors} &
  k\text{-string tensions and excited flux tubes}\\
\text{large-}N_c\text{ expansion} &
  M_{\rm glueball}=O(1),\quad \Gamma_{\rm glueball}=O(N_c^{-2})
\end{array}
\]
The bosonic theory supplies the target spectral questions: a mass gap,
isolated glueball poles, stable flux sectors, and an effective string
description in the long-string regime.

\subsection[Pure N=1 Yang--Mills]{Pure \(\mathcal N=1\) Yang--Mills}

Pure \(\mathcal N=1\) Yang--Mills adds protected chiral branch data.  The
local coordinate
\[
  S=-\frac{1}{32\pi^2}\operatorname{tr}W^\alpha W_\alpha
\]
has branch values
\[
  \langle S\rangle_k=C\Lambda_h^3\ee^{2\pi\ii k/N_c}.
\]
The family ledger separates three outputs:
\[
\begin{array}{c|c}
\text{output} & \text{source}\\
\hline
N_c\text{ branch phases} & \mathbb Z_{2N_c}\text{ anomaly and chiral }F\text{-term}\\
T_{ij}=2|W_i-W_j| & \mathcal N=1\text{ BPS wall algebra}\\
\text{glueball/gluinoball masses} & \text{spectral measure }\rho_{ij}^{(k)}
\end{array}
\]
The first two are protected after the stated hypotheses of
Chapter~\ref{ch:susy-four-dimensional-n1-gauge-dynamics}.  The third is a
massive spectral problem involving the scalar, pseudoscalar, and fermionic
components of gauge-invariant superfields.  Supersymmetry organizes isolated
massive states into multiplets when the vacuum preserves the supercharges,
but it does not supply their numerical masses.

\begin{definition}[Pure-SYM channel source ledger]
\label{def:pure-sym-channel-source-ledger}
Fix a massive branch \(k\) of pure \(\mathcal N=1\) \(SU(N_c)\) Yang--Mills
and the vacuum spectral datum of
Definition~\ref{def:pure-sym-vacuum-spectral-datum}.  Let
\[
  S=-\frac{1}{32\pi^2}\operatorname{tr}W^\alpha W_\alpha
\]
be the chiral coordinate used in
Section~\ref{sec:pure-sym-protected-spectral-data}.  The elementary source
ledger attached to this superfield is:
\begin{center}
\small
\begin{tabular}{p{0.20\textwidth}|p{0.37\textwidth}|p{0.32\textwidth}}
\textbf{channel} & \textbf{source family} & \textbf{spectral role}\\
\hline
bosonic even scalar &
  \(\operatorname{Re} S|_{\theta=0}\), \(\operatorname{Re}D^2S|\), and
  same-quantum-number descendants &
  overlaps with \(0^{++}\)-type atoms after operator mixing is fixed\\
bosonic odd scalar &
  \(\operatorname{Im} S|_{\theta=0}\), \(\operatorname{Im}D^2S|\), and
  same-quantum-number descendants &
  overlaps with \(0^{-+}\)-type atoms after operator mixing is fixed\\
fermionic &
  \(D_\alpha S|\), \(\overline D_{\dot\alpha}\overline S|\), and
  derivative descendants &
  overlaps with spin-\(\frac12\) atoms in the same massive multiplets\\
\end{tabular}
\end{center}
Here the vertical bar denotes restriction to \(\theta=\overline\theta=0\) after
the chosen superspace operator has been renormalized.  The table is a source
ledger, not a list of particles: the spectral measure may contain several
isolated atoms and continuum components in any row, and the same atom can have
nonzero residue for several sources with the same exact quantum numbers.
Center-charged Wilson or magnetic lines belong to a separate line-sector
ledger; they define flux-string sectors rather than chiral-particle masses.
\end{definition}

\begin{proposition}[Mass degeneracy inside an isolated pure-SYM multiplet]
\label{prop:pure-sym-isolated-multiplet-mass-degeneracy}
Assume the vacuum branch \(k\) carries unbroken \(\mathcal N=1\) supercharges
\(Q_\alpha,\overline Q_{\dot\alpha}\) on the Hilbert space of
Definition~\ref{def:pure-sym-vacuum-spectral-datum}.  Let \(E_k(\Sigma_m)\)
be the spectral projection onto an isolated positive-mass shell
\(\Sigma_m=\{p:p^2=-m^2,\ p^0>0\}\), separated from the rest of the joint
translation spectrum in a neighborhood invariant under the little group.
Then the range of \(E_k(\Sigma_m)\) is invariant under the supercharges.  In
particular, whenever a bosonic or fermionic source in
Definition~\ref{def:pure-sym-channel-source-ledger} has nonzero overlap with
this atom, the complete supermultiplet generated by the supercharges has the
same mass \(m\).
\end{proposition}

\begin{proof}
The supersymmetry algebra gives
\[
  [P_\mu,Q_\alpha]=[P_\mu,\overline Q_{\dot\alpha}]=0.
\]
Therefore the mass operator
\[
  M_k^2=-\eta^{\mu\nu}P_\mu^{(k)}P_\nu^{(k)}
\]
commutes with the supercharges on their common invariant domain.  Functional
calculus then gives
\[
  Q_\alpha E_k(B)=E_k(B)Q_\alpha,\qquad
  \overline Q_{\dot\alpha}E_k(B)=E_k(B)\overline Q_{\dot\alpha}
\]
for every Borel set \(B\) of the joint spectrum for which the products are
defined by closure; applying this to the isolated shell gives invariance of
\(\operatorname{Ran}E_k(\Sigma_m)\).  Acting with the supercharges changes
spin and fermion number but not \(P^2\), so all nonzero states in the generated
irreducible massive supermultiplet have the same mass.  The statement says
nothing about the numerical value of \(m\) or about whether the atom exists;
those are spectral questions in the chosen nonperturbative construction.
\end{proof}

\begin{definition}[Regulated pure-SYM spectral extraction datum]
\label{def:regulated-pure-sym-spectral-extraction-datum}
A regulated extraction of the pure \(\mathcal N=1\) SYM spectrum consists of:
\begin{enumerate}[label=(\roman*)]
\item a gauge-field regulator and an adjoint Majorana-fermion regulator with
  a specified Pfaffian or determinant prescription;
\item a parameter \(m_{\lambda,\mathrm{reg}}\) whose tuned value defines the
  massless-gluino trajectory, together with the Ward identity or equivalent
  renormalized condition used to tune it;
\item a finite-volume transfer matrix or Euclidean reflection-positive
  substitute sufficient to define spectral estimates;
\item source matrices in the exact finite-regulator symmetry channels, with
  the continuum \(J^{PC}\) labels assigned only after the rotation and
  discrete-symmetry restoration limits have been taken;
\item a limit prescription \(a\to0\), \(L\to\infty\), and
  \(m_{\lambda,\mathrm{reg}}\to m_{\lambda,\ast}\), including the order of
  limits and the systematic error controls used in the extrapolation.
\end{enumerate}
This datum is the supersymmetric analogue of the pure-glue datum in
Section~\ref{sec:lattice-ym-glueballs-flux-observables}.  It is not replaced
by the Veneziano--Yankielowicz superpotential: the latter supplies protected
branch data, while this datum is what can define numerical glueball,
gluinoball, and mixed supermultiplet masses.
\end{definition}

\subsection[Pure N=2 and Softly Broken N=2]{Pure \(\mathcal N=2\) and Softly Broken \(\mathcal N=2\)}

Pure \(\mathcal N=2\) Yang--Mills has Coulomb-branch coordinate \(u\), period
vector \((a_D,a)\), and BPS central charge
\[
  Z_\gamma=n_ea+n_ma_D .
\]
At a monopole singularity, a soft \(\mathcal N=1\) deformation
\[
  W_{\mathrm{loc},\mu}
  =
  \sqrt2\,A_DM\widetilde M+\mu\,U(A_D)
\]
gives
\[
  A_D=0,\qquad
  M\widetilde M=-\frac{\mu}{\sqrt2}U'(0).
\]
The dual \(U(1)_D\) is Higgsed, and the local Abelian model supports flux
tubes carrying the original electric flux.  The Seiberg--Witten row therefore
contains a controlled Abelian confinement mechanism near the singularity.
The detailed non-Abelian \(k\)-string spectrum away from the local Abelian
regime is a separate spectral problem.

\begin{controlledapproximation}[Radial BPS string profile near a monopole point]
\label{ca:sw-radial-bps-string-profile}
Use the local Abelian-Higgs normalization of
\eqref{eq:sw-local-abelian-higgs-energy}, take \(n>0\), and choose the
orientation for which
\[
  (D_1+\ii D_2)q=0,\qquad B_D=e_D^2(|\xi|-|q|^2).
  \label{eq:family-sw-bps-vortex-equations}
\]
With polar coordinate \((r,\varphi)\) on the plane transverse to the string,
set
\[
  q(r,\varphi)=|\xi|^{1/2}f(R)\ee^{\ii n\varphi},
  \qquad
  A_D=n\,a(R)\,\dd\varphi,\qquad
  R=e_D|\xi|^{1/2}r .
  \label{eq:family-sw-vortex-ansatz}
\]
Finite energy and regularity impose
\[
  f(0)=0,\qquad a(0)=0,\qquad
  f(\infty)=1,\qquad a(\infty)=1.
  \label{eq:family-sw-vortex-boundary-conditions}
\]
The first-order equations reduce to
\[
  \frac{\dd f}{\dd R}=\frac{n}{R}(1-a)f,
  \qquad
  \frac{\dd a}{\dd R}=\frac{R}{n}(1-f^2).
  \label{eq:family-sw-vortex-radial-ode}
\]
The flux and BPS tension are
\[
  \int_{\mathbb R^2}F_D=2\pi n,\qquad
  T_n=2\pi n|\xi|.
  \label{eq:family-sw-vortex-flux-tension}
\]
\end{controlledapproximation}

\begin{proof}
For the ansatz \eqref{eq:family-sw-vortex-ansatz},
\[
  D_rq=|\xi|^{1/2}e_D|\xi|^{1/2}f'(R)\ee^{\ii n\varphi},
  \qquad
  D_\varphi q=\ii n(1-a)|\xi|^{1/2}f(R)\ee^{\ii n\varphi}.
\]
Since \(D_1+\ii D_2=\ee^{\ii\varphi}(D_r+\ii r^{-1}D_\varphi)\), the first BPS
equation gives
\[
  e_D|\xi|^{1/2}f'(R)-\frac{n}{r}(1-a)f(R)=0,
\]
which is the first equation in \eqref{eq:family-sw-vortex-radial-ode}.  The
curvature of the one-form \(A_D=n\,a(R)\dd\varphi\) is
\[
  F_D=n\,a'(R)\dd R\wedge\dd\varphi,
  \qquad
  B_D=\frac{n}{r}e_D|\xi|^{1/2}a'(R).
\]
Substituting \(r=R/(e_D|\xi|^{1/2})\) into
\(B_D=e_D^2|\xi|(1-f^2)\) gives the second radial equation.  Finally,
\[
  \int_{\mathbb R^2}F_D
  =
  \int_0^{2\pi}\!\dd\varphi\int_0^\infty n\,a'(R)\dd R
  =
  2\pi n\,[a(\infty)-a(0)]
  =
  2\pi n .
\]
The tension formula is then the saturated bound
\eqref{eq:sw-local-vortex-bound} with positive orientation.  The derivation is
local in the Seiberg--Witten patch; it does not assert that a full non-Abelian
\(k\)-string spectrum is determined by the Abelian profile.
\end{proof}

\subsection[Planar N=4]{Planar \(\mathcal N=4\)}

Planar \(\mathcal N=4\) Yang--Mills is conformal, so the family ledger uses
scaling dimensions and defect energies instead of massive glueball poles.  A
local-operator row is the planar spin-chain/QSC spectral problem of
Chapters~\ref{ch:planar-n4-spectral-problem-spin-chains}--\ref{ch:planar-n4-quantum-spectral-curve-hexagon}.  A line-defect row is the
Wilson-line flux-tube datum of
Definition~\ref{def:planar-n4-wilson-line-flux-tube-datum}.  Its vacuum
energy is the cusp anomalous dimension, and its excitations are the
pentagon-OPE or mirror-TBA flux-tube excitations in an integrability
description.

The \(\mathcal N=4\) row is therefore the most exactly soluble spectral row
currently available inside four-dimensional interacting QFT, while its
proof boundary is also sharp: the planar integrability construction is an
exact spectral framework under stated integrability hypotheses, and a
derivation from a nonperturbative four-dimensional regulator remains open.

\section{Soft Breaking Toward Bosonic Yang--Mills}
\label{sec:susy-ym-family-soft-breaking-to-bosonic}

The deformation from pure \(\mathcal N=1\) Yang--Mills toward bosonic
Yang--Mills is naturally expressed by a spurion source for the chiral
coordinate \(S\):
\[
  \Delta\mathcal L
  =
  \int\dd^2\theta\,X_\lambda(\theta)S+\mathrm{c.c.},
  \qquad
  X_\lambda(\theta)=\theta^2m_\lambda .
  \label{eq:gaugino-mass-spurion-source}
\]
The component deformation is
\[
  \Delta\mathcal L_{\rm comp}
  =
  m_\lambda\,S|_{\theta=0}
  +
  \overline m_\lambda\,\overline S|_{\overline\theta=0}.
  \label{eq:gaugino-mass-component-source}
\]
At small \(|m_\lambda|\), the leading branch energy splitting is obtained
from the chiral condensate:
\[
  \Delta E_k
  =
  -2\,\operatorname{Re}\!\left(
    m_\lambda\,C\Lambda_h^3\ee^{2\pi\ii k/N_c}
  \right)
  +O(|m_\lambda|^2).
  \label{eq:gaugino-mass-branch-splitting}
\]
This formula is a first-order source-response statement in the massive
\(\mathcal N=1\) theory.  Increasing \(|m_\lambda|\) until the gaugino
decouples is a deformation path whose spectral continuity must be examined
with nonchiral observables: glueball masses, string tensions, domain-wall
fates, and line-operator sectors.  The endpoint comparison with bosonic
Yang--Mills therefore belongs to the spectral ledger of
Definition~\ref{def:susy-ym-family-observable-ledger}, not only to the
chiral branch ledger.

\begin{controlledapproximation}[Small soft mass branch selection]
\label{ca:small-soft-mass-branch-selection}
Assume pure \(\mathcal N=1\) \(SU(N_c)\) Yang--Mills has \(N_c\) massive
branches with condensates
\[
  S_k=C\Lambda_h^3\ee^{2\pi\ii k/N_c}
\]
and that the soft source \eqref{eq:gaugino-mass-spurion-source} can be
treated in first-order perturbation theory in a finite volume before the
thermodynamic limit is taken.  Then the branch selected at first order is the
branch maximizing
\[
  \operatorname{Re}(m_\lambda S_k).
\]
\end{controlledapproximation}

\begin{proof}
The first-order shift of the vacuum energy density under the component source
\(m_\lambda S+\overline m_\lambda\overline S\) is minus the expectation value
of the source term:
\[
  \Delta E_k
  =
  -
  \left(m_\lambda S_k+\overline m_\lambda\,\overline S_k\right)
  +O(|m_\lambda|^2)
  =
  -2\operatorname{Re}(m_\lambda S_k)+O(|m_\lambda|^2).
\]
At first order, minimizing \(\Delta E_k\) is the same as maximizing
\(\operatorname{Re}(m_\lambda S_k)\).
\end{proof}

\section{Status Ledger}
\label{sec:susy-ym-family-status-ledger}

The present chapter closes the elementary comparison map.  It records the
following theorem boundaries:
\[
\begin{array}{c|c}
\text{statement} & \text{status in this monograph}\\
\hline
\mathcal N=4\to\mathcal N=2,\mathcal N=1
\text{ scale matching} & \text{derived holomorphic threshold matching}\\
\mathcal N=1\text{ branches and BPS walls} &
  \text{derived from chiral }F\text{-term hypotheses}\\
\mathcal N=2\text{ monopole condensation} &
  \text{derived in local Abelian patch}\\
\mathcal N=2\text{ vortex tension} &
  \text{BPS Abelian-Higgs controlled approximation}\\
\mathcal N=4\text{ Wilson-line flux tube} &
  \text{defect spectral datum plus integrability hypotheses}\\
\text{bosonic glueball and }k\text{-string spectra} &
  \text{lattice or Hilbert-space spectral problem}
\end{array}
\]
The frontier problem is to construct a regulator-level family in which enough
of the ladder is realized simultaneously to compare the spectral measures,
line-defect Hilbert spaces, and flux-sector transfer matrices directly.

\begin{openproblem}[Regulator-level supersymmetry-breaking spectral bridge]
\label{op:regulator-level-susy-breaking-spectral-bridge}
Construct a regulator and continuum-limiting framework for the deformation
from supersymmetric Yang--Mills theories to bosonic Yang--Mills that controls
all of the following in the same topology: gauge-invariant local spectral
measures, center-flux transfer matrices, line-operator charge lattices,
domain-wall sectors, and protected chiral coordinates.  Prove, disprove, or
replace the commonly assumed continuity of glueball and flux-string spectra
along such deformation paths by a theorem with stated hypotheses.
\end{openproblem}
